% ===== ch14 : 原文 main.tex 第 5024--5545 行（附录 A Eilenberg-Zilber 收缩（后，含 align* 起））=====
% 本文件是原文 LaTeX 源码的逐行分片，仅供翻译时对照；最终译文写到 src/chunks/ch14.tex。
\begin{align*}
\begin{split}
&
\check{C}_{G}^p
\Big(
\mathfrak{U},
\mathcal{H}om_{\mathcal{O}_X}^q
\big(
(\mathcal{F}^{\bullet},\theta^{\bullet}),
(\mathcal{F}'^{\bullet},\theta'^{\bullet})
\big)
\Big)
\cong
\prod_{i_0<\cdots<i_p}
\bigoplus_{i\in\mathbb{Z}}
\Hom_{\VB_{G}(U_{i_0\cdots i_p})}
% ---- 原文行 5040 ----
\big(
(\mathcal{F}^i,\theta^i)|_{U_{i_0\cdots i_p}}),
(\mathcal{F}'^i,\theta'^{i+q}|_{U_{i_0\cdots i_p}})
\big),
\\[2mm]
&
\check{C}_{H}^r
\Big(
\mathfrak{V},
\mathcal{H}om_{\mathcal{O}_Y}^s
\big(
(\mathcal{G}^{\bullet},\eta^{\bullet}),
(\mathcal{G}'^{\bullet},\eta'^{\bullet})
\big)
\Big)
\cong
\prod_{j_0<\cdots<j_r}
\bigoplus_{j\in\mathbb{Z}}
\Hom_{\VB_{H}(V_{j_0\cdots j_p})}
\big(
% ---- 原文行 5060 ----
(\mathcal{G}^j,\eta^j)|_{V_{j_0\cdots j_r}}),
(\mathcal{G}'^j,\eta'^{j+s}|_{V_{j_0\cdots j_r}})
\big),
\\[2mm]
&
\check{C}_{G\times H}^t
\Big(
\mathfrak{U}\times\mathfrak{V},
\mathcal{H}om_{\mathcal{O}_{X\times Y}}^u
\big(
(\mathcal{F}^{\bullet},\theta^{\bullet})
\boxtimes
(\mathcal{G}^{\bullet},\eta^{\bullet}),
(\mathcal{F}'^{\bullet},\theta'^{\bullet})
\boxtimes
(\mathcal{G}'^{\bullet},\eta'^{\bullet})
\big)
\Big)
\\[2mm]
&
% ---- 原文行 5080 ----
\hspace{-5mm}
\cong
\prod_{\alpha_0<\cdots<\alpha_t}
\bigoplus_{k\in\mathbb{Z}}
\bigoplus_{\substack{
i+j=k
\\[1mm]
q+s=u
}}
\Hom_{\VB_{G\times H}(W_{\alpha_0\cdots \alpha_t})}
\Big(
\big(
(\mathcal{F}^i,\theta^i)
\boxtimes
(\mathcal{G}^j,\eta^j)
\big)|_{W_{\alpha_0\cdots \alpha_t}},
\big(
(\mathcal{F}'^{i+q},\theta'^{i+q})
\boxtimes
(\mathcal{G}'^{j+s},\eta'^{j+s})
% ---- 原文行 5100 ----
\big)|_{W_{\alpha_0\cdots \alpha_t}})
\Big)
\end{split}
\end{align*}
We denote a homogeneous element in
\begin{center}
\resizebox{\textwidth}{!}{$\begin{aligned}
C_{\otimes}^{n}
&=
\bigoplus_{l+m=n}
\Big(
\hom_{\Perf_{G}(X;\mathfrak{U})_{\dg}}^l
\big(
(\mathcal{F}^{\bullet},\theta^{\bullet}),
(\mathcal{F}'^{\bullet},\theta'^{\bullet})
\big)
\underset{\field}{\otimes}
\hom_{\Perf_{H}(Y;\mathfrak{V})_{\dg}}^m
\big(
(\mathcal{G}^{\bullet},\eta^{\bullet}),
% ---- 原文行 5120 ----
(\mathcal{G}'^{\bullet},\eta'^{\bullet})
\big)
\Big)
\\[2mm]
&
\hspace{-5mm}
\cong
\bigoplus_{l+m=n}
\Bigg(
\bigg(
\bigoplus_{p+q=l}
\check{C}_{G}^p
\Big(
\mathfrak{U},
\mathcal{H}om_{\mathcal{O}_X}^q
\big(
(\mathcal{F}^{\bullet},\theta^{\bullet}),
(\mathcal{F}'^{\bullet},\theta'^{\bullet})
\big)
\Big)
% ---- 原文行 5140 ----
\bigg)
\underset{\field}{\otimes}
\bigg(
\bigoplus_{r+s=m}
\check{C}_{H}^r
\Big(
\mathfrak{V},
\mathcal{H}om_{\mathcal{O}_Y}^s
\big(
(\mathcal{G}^{\bullet},\eta^{\bullet}),
(\mathcal{G}'^{\bullet},\eta'^{\bullet})
\big)
\Big)
\bigg)
\Bigg)
\\[2mm]
&
\hspace{-5mm}
\cong
\bigoplus_{l+m=n}
% ---- 原文行 5160 ----
\bigoplus_{\substack{
p+q=l
\\[1mm]
r+s=m
}}
\prod_{\substack{
i_0<\cdots<i_p
\\[0mm]
j_0<\cdots<j_r
}}
\bigoplus_{\substack{
i\in\mathbb{Z}
\\[1mm]
j\in\mathbb{Z}
}}
\Big(
\Hom_{\VB_{G}(U_{i_0\cdots i_p})}
\big(
(\mathcal{F}^i,\theta^i)|_{U_{i_0\cdots i_p}},
(\mathcal{F}'^{i+q},\theta'^{i+q})|_{U_{i_0\cdots i_p}}
% ---- 原文行 5180 ----
\big)
\\
&\phantom{
\cong
\bigoplus_{l+m=n}
\bigoplus_{\substack{
p+q=l
\\[1mm]
r+s=m
}}
\prod_{\substack{
i_0<\cdots<i_p
\\[0mm]
j_0<\cdots<j_r
}}
\bigoplus_{\substack{
i\in\mathbb{Z}
\\[1mm]
j\in\mathbb{Z}
}}
% ---- 原文行 5200 ----
\Big(
}
\underset{\field}{\otimes}
\Hom_{\VB_{H}(V_{j_0\cdots j_r})}
\big(
(\mathcal{G}^j,\eta^j)|_{V_{j_0\cdots j_r}},
(\mathcal{G}'^{j+s},\eta'^{j+s})|_{V_{j_0\cdots j_s}}
\big)
\Big)
\end{aligned}$}
\end{center}
of multi-degree
$(p,q,r,s)$
by
$$
f\otimes g
=
\Big(
f_{i_0\cdots i_p}^i
\otimes
% ---- 原文行 5220 ----
g_{j_0\cdots j_r}^j
\Big)
_{\substack{
i,j\in\mathbb{Z}
\\
i_0<\cdots<i_p
\\
j_0<\cdots<j_r
}}.
$$


The Eilenberg--Zilber maps
$S$, $T$ and $H$
in
(\ref{eqn:EZMaps})
induce
cochain maps
$S^*$, $T^*$ and$H^*$
between $C_{\otimes}^{\bullet}$ and $C_{\times}^{\bullet}$:
% ---- 原文行 5240 ----



$$
\begin{tikzcd}[arrow style=tikz,>=stealth,row sep=2.5em,column sep=3.5em]
C_{\otimes}^{\bullet}
  \arrow[r, shift left, <-, "S^*"]
%  \arrow[loop left, <-, "H"]
&
C_{\times}^{\bullet}
  \arrow[l, shift left, <-, "T^*"]
  \arrow[loop right, "\mathbb{H}^{\ast}"]
\end{tikzcd}
$$


\begin{dfn}

(1)
The chain map 
% ---- 原文行 5260 ----
$
T:
C_{\bullet}(I\times J)
\rightarrow
C_{\bullet}(I)\otimes C_{\bullet}(J)
$
induces a map between cochain complexes
\begin{align*}
T^*:
C_{\otimes}^n
&\rightarrow
C_{\times}^n
\\
\quad
f\otimes g
=
\Big(
f_{i_0\cdots i_p}^i
\otimes
g_{j_0\cdots j_r}^j
% ---- 原文行 5280 ----
\Big)
_{\substack{
i,j\in\mathbb{Z}
\\
i_0<\cdots<i_p
\\
j_0<\cdots<j_r
}}
&\mapsto
T^*(f\otimes g)
:=
\Big(
T^*(f\otimes g)_{\alpha_0\cdots\alpha_t}^{i,j}
\Big)
_{\substack{
i,j\in\mathbb{Z}
\\
i_0<\cdots<i_p
\\
j_0<\cdots<j_r
% ---- 原文行 5300 ----
}},
\end{align*}
whose components are defined by
$$
T^*(f\otimes g)_{\alpha_0\cdots\alpha_t}^{i,j}
:=
\sum_{\substack{
i_0<\cdots<i_p
\\
j_0<\cdots<j_r
}}
(-1)^{qr+qs+si}
T_{\alpha_0\cdots\alpha_t}^{i_0\cdots i_p;j_0\cdots j_r}
\big(
f_{i_0\cdots i_p}^i
\boxtimes
g_{j_0\cdots j_r}^j
\big)\big|_{W_{\alpha_0\cdots\alpha_t}},
$$
where the coefficients
% ---- 原文行 5320 ----
$
T_{\alpha_0\cdots\alpha_t}^{i_0\cdots i_p;j_0\cdots j_r}
\in\mathbb{Z}
$
are determined by the formula
$$
T\big((\alpha_0,\dots,\alpha_t)\big)
=
\sum_{\substack{
i_0<\cdots<i_p
\\
j_0<\cdots<j_r
}}
T_{\alpha_0\cdots\alpha_t}^{i_0\cdots i_p;j_0\cdots j_r}
(i_0,\dots,i_p)\otimes(j_0,\dots,j_r).
$$

It follows from the definition of $T$
(Definition \ref{defn:EZMaps}.(2))
that
% ---- 原文行 5340 ----
$$
T_{\alpha_0\cdots\alpha_t}^{i_0\cdots i_p;j_0\cdots j_r}
=
\begin{cases}
1
&
\text{if }
(\alpha_0,\dots,\alpha_t)
=
\big((i_0,*),\dots,(i_{p-1},*),(i_p,j_0),(*,j_1),\dots,(*,j_r)\big),
\\
0
&
\text{otherwise},
\end{cases}
$$
and therefore we have
$$
T^*(f\otimes g)_{\alpha_0\cdots\alpha_t}^{i,j}
=
% ---- 原文行 5360 ----
(-1)^{qr+qs+si}
\big(
f_{i_0\cdots i_p}^i
\boxtimes
g_{j_p\cdots j_t}^j
\big)\big|_{W_{\alpha_0\cdots\alpha_t}}
$$
where
$(\alpha_0,\dots,\alpha_t)=\big((i_0,j_0),\dots,(i_t,j_t)\big)$.

(2)
The chain map
$
S:
C_{\bullet}(I)\otimes C_{\bullet}(J)
\rightarrow
C_{\bullet}(I\times J)
$
induces
a map between cochain complexes
% ---- 原文行 5380 ----
\begin{align*}
S^*:
C_{\times}^n
&\rightarrow
C_{\otimes}^n
\\
\quad
h
=
\Big(
h_{\alpha_0\cdots\alpha_t}^{i,j}
\Big)
_{\substack{
i,j\in\mathbb{Z}
\\
\alpha_0<\cdots<\alpha_t
}}
&\mapsto
S^*h
:=
% ---- 原文行 5400 ----
\Big(
(S^*h)_{i_0\cdots i_p;j_0\cdots j_r}^{i,j}
\Big)
_{\substack{
i,j\in\mathbb{Z}\phantom{,\ p+r=t}
\\
p,r\in\mathbb{Z},\ p+r=t
\\
i_0<\cdots<i_p\quad\
\\
j_0<\cdots<j_r\quad\
}},
\end{align*}
whose components are defined by
$$
(S^*h)_{i_0\cdots i_p;j_0\cdots j_r}^{i,j}
:=
\sum_{\substack{
\alpha_0\cdots\alpha_t
}}
% ---- 原文行 5420 ----
(-1)^{qr+qs+si}
S_{i_0\cdots i_p;j_0\cdots j_r}^{\alpha_0\cdots\alpha_t}
\boxtimes^{-1}\big(h_{\alpha_0\cdots\alpha_t}^{i,j}\big),
$$
where the coefficients
$
S_{i_0\cdots i_p;j_0\cdots j_r}^{\alpha_0\cdots\alpha_t}
\in\mathbb{Z}
$
are determined by the formula
$$
S\big((i_0,\dots,i_p)\otimes(j_0,\dots,j_r)\big)
=
\sum_{\substack{
\alpha_0<\cdots<\alpha_t
}}
S_{i_0\cdots i_p;j_0\cdots j_r}^{\alpha_0\cdots\alpha_t}
(\alpha_0,\dots,\alpha_t)
$$
and\\
% ---- 原文行 5440 ----
\adjustbox{max width=\textwidth}{$\displaystyle
\boxtimes:
\Hom_{\VB_G(U_{i_0\cdots i_p})}
\big(
(\mathcal{F}^i,\theta^i)|_{U_{i_0\cdots i_p}},
(\mathcal{F}'^{i+q},\theta'^{i+q})|_{U_{i_0\cdots i_p}}
\big)
\underset{\field}{\otimes}
\Hom_{\VB_H(V_{j_0\cdots j_r})}
\big(
(\mathcal{G}^j,\eta^j)|_{V_{j_0\cdots j_r}},
(\mathcal{G}'^{j+s},\eta'^{j+s})|_{V_{j_0\cdots j_r}}
\big)
$}\\
\adjustbox{max width=\textwidth}{
$\displaystyle
\quad
\xrightarrow{\cong}
\Hom_{\VB_{G\times H}(U_{i_0\cdots i_p}\times V_{j_0\cdots j_r})}
\Big(
% ---- 原文行 5460 ----
\big(
(\mathcal{F}^i,\theta^i)
\boxtimes
(\mathcal{G}^j,\eta^j)
\big)
\big|_{U_{i_0\cdots i_p}\times V_{j_0\cdots j_r}},
\big(
(\mathcal{F}'^{i+q},\theta'^{i+q})
\boxtimes
(\mathcal{G}'^{j+s},\eta'^{j+s}
\big)
\big|_{U_{i_0\cdots i_p}\times V_{j_0\cdots j_r}}
\Big)
$}\\
is the isomorphism given in Proposition \ref{prop:BoxTensorFunctorFullyFaithful}.

It follows from the definition of $S$
(Definition \ref{defn:EZMaps}.(1))
that
$$
% ---- 原文行 5480 ----
(S^*h)_{i_0\cdots i_p;j_0\cdots j_r}^{i,j}
=
(-1)^{qr+qs+si}
\boxtimes^{-1}
\bigg(
\sum_{
0=l_0\le l_1\le\cdots l_{p+1}=r
}
(-1)^{l_1+\cdots+l_p}
h_{
(i_0,j_0),\dots,(i_0,j_{l_0}),
(i_1,j_{l_1}),\dots,(i_1,j_{l_2}),
\dots,
(i_p,j_{l_p}),\dots,(i_p,j_r)
}
^{i,j}
\bigg).
$$


% ---- 原文行 5500 ----
\noindent
(3)
The chain homotopy
$
\mathbb{H}:
C_{\bullet}(I\times J)
\rightarrow
C_{\bullet}(I\times J)
$
induces a map between cochain complexes
$$
\mathbb{H}^{\ast}\colon C_{\times}^{\bullet}\to C_{\times}^{\bullet -1}
$$where for each $f\coloneqq\left(f_{\beta_0\dotsb\beta_{t+1}^k}\right)_
{
    \begin{subarray}{c}
        k\in\mathbb{Z}\\ 
        \beta_0<\dotsb<\beta_{t+1}
    \end{subarray}
    }$, we have$$
(\mathbb{H}^{\ast} f)\coloneqq\left((\mathbb{H}^{\ast} f)^k_{\alpha_0\dotsb\alpha_t}\coloneqq\sum_{\beta_0<\dotsb<\beta_{t+1}}\left. \mathbb{H}_{\alpha_0\dotsb\alpha_t}^{\beta_0\dotsb\beta_{t+1}} f_{\beta_0\dotsb\beta_{t+1}}^k\right|_{W_{\alpha_0\dotsb\alpha_t}}\right)
% ---- 原文行 5520 ----
_{
    \begin{subarray}{c}
        k\in\mathbb{Z}\\ 
        \alpha_0<\dotsb<\alpha_t
    \end{subarray}
}
$$where we write \[
    \mathbb{H}((\alpha_0,\dotsb,\alpha_t))\eqqcolon\sum_{\beta_0<\dotsb<\beta_{t+1}}\mathbb{H}_{\alpha_0\dotsb\alpha_t}^{\beta_0\dotsb\beta_{t+1}} (\beta_0,\dotsb ,\beta_{t+1}).
\]
\end{dfn}
Now we can verify that
\begin{prp}\label{prp:STHd-relation}
    \begin{itemize}
        \item Write $\delta$ for the \v{C}ech differential, then $\delta\mathbb{H}^{\ast} +\mathbb{H}^{\ast} \delta=S^{\ast} T^{\ast} -\id$.
        \item $S^{\ast}$ preserves compositions, i.e. $S^{\ast} (f\comp g)=S^{\ast} f\comp S^{\ast} g$.
        \item Write $d_{\mathcal{H}^{\bullet}}$ to be the differential for the internal hom complex, then $d_{\mathcal{H}^{\bullet}}\mathbb{H}^{\ast} =\mathbb{H}^{\ast}d_{\mathcal{H}^{\bullet}}$.
        \item $dT^{\ast} =T^{\ast} d$.
        \item $T^{\ast} S^{\ast} =\id$.
    \end{itemize}
\end{prp}
% ---- 原文行 5540 ----
\begin{proof}
    These are all direct verifications, using the definition and relations in Proposition \ref{prp:STHd-relation}.
\end{proof}


% \section{K\"unneth Bimodules}
